\documentclass[11pt,a4paper]{article}

\usepackage[UTF8]{ctex}
\usepackage{amsmath,amssymb}
\usepackage{geometry}
\usepackage{enumitem}
\usepackage{booktabs}
\usepackage{xcolor}
\usepackage{hyperref}

\geometry{margin=2.2cm}

\newcommand{\key}[1]{\textcolor{blue}{\textbf{#1}}}
\newcommand{\lec}[1]{\textcolor{gray}{[lec-#1]}}
\newcommand{\ans}{\par\noindent\textcolor{blue}{\textbf{解：}}}

\title{CS 2053M 计算机系统 III\\24--25 春夏学期 期末回忆卷}
\author{整理自 CC98 论坛 + 考场抄写，题目不全，仅供参考}
\date{}

\begin{document}

\maketitle

\section*{一、选择题（共 15 道）}

\begin{enumerate}[leftmargin=*,itemsep=8pt]

\item 在没有前递（forwarding）、分支预测，且指令内存和数据内存分开的五级流水线中运
  行下列指令。在 I5 之前，发生了哪种或哪些 Hazard？

  \begin{verbatim}
  I1: LW   x1, 0(x2)
  I2: ADD  x3, x1, x4
  I3: SUB  x5, x3, x6
  I4: BEQ  x3, x0, L1
  I5: ADDI x7, x7, 1
  \end{verbatim}

  \ans RAW hazard：I1 $\to$ I2（x1），I2 $\to$ I3（x3），I2 $\to$ I4（x3）；
  Control hazard：I4（BEQ）。无 forwarding 时所有 RAW 都需要 stall；
  无分支预测则 BEQ 须等 EX 阶段算出结果，fetch 阻塞。 \lec{02}

\item 下列哪种硬件结构\textbf{不会} typically 和动态调度（dynamic scheduling）联
  系在一起？

  \begin{enumerate}[label=\Alph*.]
    \item Branch Delay Slots
    \item Reservation Station
    \item ROB (Reorder Buffer)
    \item Register Rename Table
  \end{enumerate}

  \ans A。Branch Delay Slot 是静态调度技术（编译器填充），Reservation Station、
  ROB、Register Rename Table 三者是 Tomasulo 动态调度的核心硬件。 \lec{03}

\item 有关记分牌（Scoreboard）算法和 Tomasulo 方法的区别，以下说法正确的是？

  \ans Scoreboard 只做动态调度但不消除 WAR/WAW（靠 stall 等待），
  Tomasulo 用寄存器重命名\textbf{消除} WAR/WAW，两者都处理 RAW。
  Tomasulo 将 Scoreboard 的控制集中式改为 Reservation Station 分布式。
  且 Tomasulo 配合 ROB 支持 Hardware Speculation，Scoreboard 不支持。 \lec{03}

\item Cache 大小为 16\,KB，Block Size 为 64\,B，缓存是四路组相联的，求 set 数量。

  \ans $\text{\#blocks}=16\times1024/64=256$，
  $\text{\#sets}=256/4=64$ 个 set。 \lec{04,05}

\item 对于直接映射（DM）、全相联（FA）、组相联（SA）缓存的 Miss Rate 而言
  （假设容量和 Block Size 都一样）：

  \begin{enumerate}[label=\Alph*.]
    \item DM always 有更低的 Miss rate
    \item FA usually 有更低的 Miss rate
    \item SA always 比 DM 有更高的 Miss rate
    \item 由于 Size \& Block Size 一样，三者的 Cache Miss rate 相等
  \end{enumerate}

  \ans B。FA 无冲突缺失（conflict miss），在同等容量下通常 Miss Rate 最低。
  DM 冲突缺失最多。SA 介于两者之间。A/C 错在 ``always''；D 忽略了映射
  方式对冲突缺失的影响。 \lec{04,05}

\item 4\,block、2-way set associative Cache，block size = 1\,word，采用写回~+
  写分配策略，LRU 替换。对以下访问序列求 Cache miss 数量：

  \begin{verbatim}
  Read 0, Read 4, Read 0, Read 8, Read 4, Write 12, Write 0
  \end{verbatim}

  \ans 4 blocks, 2-way SA $\to$ 2 sets（set index = address mod 2）。

  \begin{tabular}{c|c|c|c}
  Access & Set & Hit/Miss & Set state (LRU left) \\
  \midrule
  R0 & 0 & Miss & [0,\_] \\
  R4 & 0 & Miss & [0,4] \\
  R0 & 0 & Hit  & [4,0] \\
  R8 & 0 & Miss, evict 4 & [0,8] \\
  R4 & 0 & Miss, evict 0 & [8,4] \\
  W12& 0 & Miss, evict 8 & [4,12] \\
  W0 & 0 & Miss, evict 4 & [12,0] \\
  \end{tabular}

  \textbf{总 miss = 6}（只有第 3 次 R0 命中）。 \lec{04,05}

\item 有关 Vector Processors 的概念判断：以下说法正确的是？

  \ans 向量处理器是 SIMD 架构的实现：一条向量指令同时对一组向量元素做
  相同操作。向量元素之间很少有相关性，所以向量流水线效率很高。
  向量处理器 \$\neq\$ 标量处理器（Scalar Processor），前者有向量数据表示和
  相应的向量指令集。 \lec{18}

\item 在缓存全命中情况下 CPI = 1.0，Load/Store 占总指令 50\%，Miss Rate = 2\%，
  Miss Penalty = 20 时钟周期。求在保证全部缓存命中下的 Speedup。

  \ans
  $\text{Memory stall cycles} = \text{IC} \times 0.5 \times 0.02 \times 20
  = 0.2\,\text{IC}$。

  实际 CPI：$\text{CPI}_\text{actual} = 1.0 + 0.2 = 1.2$。

  全部命中 CPI：$\text{CPI}_\text{ideal} = 1.0$。

  $\text{Speedup} = 1.2 / 1.0 = \mathbf{1.2}$。 \lec{07}

\item SCAN 算法（电梯算法）读下列 Cylinder（假设当前磁头在 50，向大号方向移动，
  请求队列：95, 180, 34, 119, 11, 123, 62, 64），求磁头移动总路程。

  \ans SCAN 从 50 出发向右到 180，再折返向左到 11：
  $50\to62\to64\to95\to119\to123\to180\to34\to11$。

  路程 $= (62-50)+(64-62)+(95-64)+(119-95)+(123-119)+(180-123)+
  (180-34)+(34-11) = 12+2+31+24+4+57+146+23 = \mathbf{299}$。 \lec{13}

\item （页表相关）某系统采用 32-bit 虚拟地址、4\,KB 页面大小，两级页表
  （每级 1024 个 PTE，每个 PTE 4\,B）。求：每级页表大小，以及
  虚拟地址 VPN 各段的位数划分。

  \ans 4\,KB 页 $\to$ offset 12\,bit。两级各 1024 项 $\to$ 每级 10\,bit。
  每级页表大小 $= 1024 \times 4 = 4\,\text{KB}$，正好一页。
  划分：{\small [VPN1 10\,bit | VPN0 10\,bit | offset 12\,bit]}。 \lec{09}

\item （VM 相关）虚拟内存和物理内存的关系以下哪项正确？
  \begin{enumerate}[label=\Alph*.]
    \item 虚拟地址空间不能大于物理地址空间
    \item 虚拟内存在物理内存不足时仍可让程序运行
    \item 每个进程共享同一张页表
    \item 虚拟内存就是交换空间（swap）
  \end{enumerate}

  \ans B。虚拟地址空间可以远超物理内存（按需调页），每个进程有独立页表，
  交换空间只是虚拟内存的实现手段之一，不是等价概念。 \lec{10}

\item RISC-V 架构下的 PTE 不包含下列哪个元素？

  \begin{enumerate}[label=\Alph*.]
    \item Logic Page Number
    \item Physical Page Frame Number
    \item Page Residence (Valid) bit
    \item Page Permissions
  \end{enumerate}

  \ans A。PTE 存的是 PPN（物理页框号）、valid/dirty/accessed 等标志位
  和 R/W/X 权限位。VPN（逻辑页号）\textbf{不在 PTE 里}——它作为 offset 在
  页表中索引 PTE，不存储在 PTE 内部。 \lec{09,11}

\item 使用 LRU 替换策略，页框初始为空。对以下引用串求 Page Fault 次数
  （假设分配 3 个页框）：

  \begin{verbatim}
  2 3 2 1 5 2 4 5 3 2 5 2
  \end{verbatim}

  \ans
  \begin{tabular}{c|c|c|c}
  访问 & 页框 & PF? & 备注 \\
  \midrule
  2 & [2,\_,\_]       & PF=1 \\
  3 & [2,3,\_]        & PF=2 \\
  2 & [2,3,\_]        & hit \\
  1 & [2,3,1]         & PF=3 \\
  5 & [5,3,1]         & PF=4, LRU=2 \\
  2 & [5,2,1]         & PF=5, LRU=3 \\
  4 & [5,2,4]         & PF=6, LRU=1 \\
  5 & [5,2,4]         & hit \\
  3 & [3,2,5]         & PF=7, LRU=4 \\
  2 & [3,2,5]         & hit \\
  5 & [3,2,5]         & hit \\
  2 & [3,2,5]         & hit \\
  \end{tabular}

  总 Page Fault = \textbf{7}。 \lec{10}

\item （分支预测/猜测执行相关）在 Hardware-Based Speculation 中，以下说法
  \textbf{错误}的是？

  \begin{enumerate}[label=\Alph*.]
    \item ROB 在 Commit 阶段才将结果写回寄存器文件
    \item 分支预测失败时，ROB 中错误路径的指令会被 flush
    \item Speculation 允许指令在分支结果确认前执行
    \item Write Result 阶段直接将结果写入寄存器文件
  \end{enumerate}

  \ans D。在 Hardware Speculation 中，Write Result 阶段将结果写入 ROB
  （而非直接写寄存器文件），Commit 阶段才从 ROB 写回寄存器文件。
  这保证了错误的推测结果不会污染体系结构状态。 \lec{03}

\item （综合判断）关于 DLP 和 TLP 的对比，以下说法正确的是？

  \begin{enumerate}[label=\Alph*.]
    \item SIMD 属于 TLP 的实现
    \item MIMD 利用 TLP，适合并行执行多个线程
    \item Flynn 分类中，SISD 包含向量处理器
    \item DLP 通过提高时钟频率来加速程序
  \end{enumerate}

  \ans B。SIMD 属于 DLP（数据级并行），不是 TLP。MIMD 利用 TLP 让
  多个独立线程在不同处理器/核上同时运行。向量处理器是 SIMD 不是 SISD。
  DLP 一次操作多个数据，不靠提频。 \lec{18,19}

\end{enumerate}

\section*{二、大题}

\subsection*{1. Hardware-Based Speculation（15'）}

设仅有\textbf{一个} Reservation Station（RS）+ ROB。
加/减法 1 周期，乘法 6 周期，除法 20 周期。
填写下表（所有指令在 ROB 中按序提交）：

\begin{verbatim}
inst          operands    issue    execute    write result    commit
div  x2,x3,x4   x3,x4
mul  x1,x5,x6   x5,x6
add  x3,x7,x8   x7,x8
mul  x1,x1,x3   x1,x3
sub  x4,x1,x5   x1,x5
\end{verbatim}

\ans 仅 1 个 RS $\to$ 指令须串行发射（前一条 WB 释放 RS 后才发下一条）。
ROB 顺序：div(1) $\to$ mul(2) $\to$ add(3) $\to$ mul(4) $\to$ sub(5)，
按序提交。

\vspace{4pt}
\begin{tabular}{c|c|c|c|c|c}
inst & operands & issue & execute & write result & commit \\
\midrule
div  & x3,x4 & 1  & 1--20 & 21 & 21 \\
mul  & x5,x6 & 22 & 22--27 & 28 & 28 \\
add  & x7,x8 & 29 & 29 & 30 & 30 \\
mul  & x1,x3 & 31 & 31--36 & 37 & 37 \\
sub  & x1,x5 & 38 & 38 & 39 & 39 \\
\end{tabular}

\vspace{4pt}
\noindent 说明：
\begin{itemize}[nosep]
  \item div issue@c1，EX 1--20，WB@c21，RS 释放；
  \item mul x1,x5,x6 issue@c22（RS 空，x5/x6 就绪），EX 22--27，WB@c28；
  \item add issue@c29，EX 1 周期，WB@c30；
  \item 第二条 mul 的 x1 来自第一条 mul（WB@c28），x3 来自 add（WB@c30），
        最早 issue@c31，EX 31--36，WB@c37；
  \item sub issue@c38（等 mul WB@c37 释放 RS），EX 1 周期，WB@c39。
\end{itemize}

\noindent 注：单 RS 是\textbf{严重瓶颈}，真实 Tomasulo 每种 FU 有独立 RS 可并行发射。
\textbf{本题令 RS 串行化是故意考察对 issue 条件的理解。} \lec{03}

\subsection*{2. MSI Cache Coherence Protocol（15'）}

图与示例与课堂 PPT 相同（三处理器 C0, C1, C2，总线 snooping）。

\begin{enumerate}[label=(\arabic*),leftmargin=*]
  \item C0: W, A11C $\leftarrow$ 012C
  \item C1: W, A108 $\leftarrow$ 0308
  \item C2: R, A108
\end{enumerate}

\ans MSI 三状态：\textbf{M}odified（脏，独占）、\textbf{S}hared（干净，可能多副本）、
\textbf{I}nvalid。

\begin{enumerate}[label=(\arabic*),nosep]
  \item \textbf{C0 写 A11C}：C0 发 BusRdX 获取 A11C 的独占副本，
        其他缓存中该行标 I。C0 拿到后写入 012C，状态 $\to$ M。
  \item \textbf{C1 写 A108}：C1 发 BusRdX，若 C0 持 M 态则需要写回
        （write-back）并标 I。C1 拿到独占副本，写入 0308，状态 $\to$ M。
  \item \textbf{C2 读 A108}：C1（M 态）snoop 到 BusRd，将数据写回内存，
        C1 $\to$ S，C2 从总线拿数据 $\to$ S（共享）。
\end{enumerate}

\noindent 关键：M 态被 snoop 命中时必须 flush 到内存再转 S；所有状态转换由
总线事务触发。 \lec{3-4?}\ \textcolor{red}{\textcolor{red}{!!!} 课件缺 MSI 内容，需自行补充}

\subsection*{3. RISC-V Sv39 虚拟内存（20'）}

RISC-V Sv39 地址翻译图（与 Lab 3 实验指导中相同）。

\begin{enumerate}[label=(\arabic*),leftmargin=*]
  \item 分别阐述 segmentation 与 paging 各自的优劣势。
  \item RISC-V 64 是否支持 segmentation？
\end{enumerate}

给定代码：
\begin{verbatim}
int* buf = NULL;
int main() {
    buf = (int*)malloc(4096 * 16);
    buf[0] = 1919;
}
\end{verbatim}

\begin{enumerate}[label=(\arabic*),leftmargin=*,resume]
  \item
  \begin{enumerate}[label=(\roman*)]
    \item \texttt{malloc} 一句是否分配了物理内存？
    \item 哪行触发了 page fault？
    \item 整段代码运行时一共分配了多少物理空间？
  \end{enumerate}
  \item
  \begin{enumerate}[label=(\roman*)]
    \item 整段代码运行过程中一共触发了几次 page fault？
    \item page fault 分别是由哪些页导致的？
  \end{enumerate}
  \item 写出虚拟地址 \texttt{0xc0000000} $\to$ 物理地址 \texttt{0xb000}
        的完整地址翻译过程（Sv39 三级页表 walk）。
\end{enumerate}

\ans

\begin{enumerate}[label=(\arabic*),leftmargin=*]

  \item
  \begin{itemize}[nosep]
    \item \textbf{Segmentation 优势}：无内部碎片，天然对应逻辑单元（代码/数据/栈），
          便于共享和保护。
    \item \textbf{Segmentation 劣势}：外部碎片严重，分配复杂（best-fit/first-fit），
          需要紧凑（compaction）。
    \item \textbf{Paging 优势}：无外部碎片，分配简单（固定大小 frame），
          支持虚拟内存（demand paging）。
    \item \textbf{Paging 劣势}：内部碎片（平均每进程 $\tfrac12$ 页），
          页表存储开销大，地址翻译多级访问内存。
  \end{itemize}

  \item RISC-V 64 \textbf{不支持} segmentation。RISC-V 的特权架构采用
        纯分页模型（Sv39/Sv48/Sv57），虚拟地址到物理地址通过多级页表翻译，
        无段寄存器或段描述符机制。 \lec{08,09}

  \item
  \begin{enumerate}[label=(\roman*)]
    \item \textbf{没有}。\texttt{malloc} 仅在堆的虚拟地址空间中\textbf{预留}
          了 $4096\times 16 = 64\,\text{KB}$（16 页），但未分配物理页框。
          物理内存采用按需调页（demand paging），实际分配发生在首次访问时。
    \item \texttt{buf[0] = 1919} 触发了 page fault。
          这是对刚 \texttt{malloc} 的虚拟地址区域的\textbf{首次写访问}。
    \item 仅分配了 \textbf{1 页 = 4\,KB} 物理内存。
          只有 \texttt{buf[0]} 所在的那一页被访问到，其他 15 页从未被触碰。
  \end{enumerate}

  \item
  \begin{enumerate}[label=(\roman*)]
    \item 共触发 \textbf{1 次} page fault（\texttt{buf[0] = 1919} 那一次）。
          程序的代码段和栈在进程启动时已被 OS 预加载（pre-fault）。
    \item 由 \texttt{buf} 指向的堆区域第 0 页导致（该页在 \texttt{malloc}
          返回的虚拟地址范围内的第一页）。
  \end{enumerate}

  \item \textbf{Sv39 三级页表翻译：} \newline
        Sv39：虚拟地址 39\,bit，页大小 4\,KB $\to$ offset 12\,bit，
        三级页表每级 9\,bit（VPN[2] : VPN[1] : VPN[0] : offset
        $= 9:9:9:12$）。

        \texttt{0xc0000000} 二进制：
        \texttt{1100\,0000\,0000\,0000\,0000\,0000\,0000\,0000}

        \begin{itemize}[nosep]
          \item VPN[2] (38:30) = \texttt{0b110000000} = \textbf{0x180}
          \item VPN[1] (29:21) = \texttt{0b000000000} = \textbf{0x000}
          \item VPN[0] (20:12) = \texttt{0b000000000} = \textbf{0x000}
          \item Offset (11:0) = \texttt{0x000}
        \end{itemize}

        翻译步骤：
        \begin{enumerate}[nosep]
          \item 读 \texttt{satp} 寄存器获取根页表 PPN；
          \item L2：读 PTE[\texttt{root + 0x180}]，得 L1 页表 PPN；
          \item L1：读 PTE[\texttt{L1 + 0x000}]，得 L0 页表 PPN；
          \item L0：读 PTE[\texttt{L0 + 0x000}]，得最终 PPN = \texttt{0xb}；
          \item PA = \texttt{0xb << 12 | 0x000} = \textbf{\texttt{0xb000}}。
        \end{enumerate}

        每次 PTE 访问都可能触发 page fault（若 PTE 无效或非叶节点），
        最坏需 3 次内存访问才完成 walk；TLB 命中则只需 1 次。 \lec{09,10,11}

\end{enumerate}

\subsection*{4. File System（20'）}

\begin{enumerate}[label=(\alph*),leftmargin=*]
  \item 采用 indexed disk block allocation（设 block size = 4\,KB，指针 4\,B），
        以下各需几个 blocks（含 index block）？
        \begin{enumerate}[label=(\roman*)]
          \item 70\,B
          \item $2^{20} + 2044$\,B
          \item $2^{20} + 2044$\,B （与 (ii) 相同，按原题保留）
        \end{enumerate}
  \item Indexed disk block allocation 能分配的最大空间（单 index block）。
  \item 假设最多使用 3 个 disk blocks 存放索引结构，求能分配的最大空间。
  \item 以下信息保存在 directory 中还是 inode 中，why？
        \begin{enumerate}[label=(\roman*)]
          \item time of creation
          \item file name
          \item access rights
        \end{enumerate}
  \item Comment the performance of \textbf{contiguous blocks}
        on reading, writing and creating files.
\end{enumerate}

\ans

\begin{enumerate}[label=(\alph*),leftmargin=*]

  \item 每 index block 可存 $4096/4 = 1024$ 个数据块指针。

        \begin{enumerate}[label=(\roman*)]
          \item 70\,B：1 数据块（70\,B $<$ 4096\,B）。
                总 = \textbf{1 index + 1 data = 2 blocks}。
          \item $2^{20}+2044 = 1\,048\,576 + 2044 = 1\,050\,620$\,B。
                $\lceil 1\,050\,620 / 4096 \rceil = 257$ 数据块。
                单 index block（1024 指针）足够。
                总 = \textbf{1 index + 257 data = 258 blocks}。
          \item 同上。（原卷回忆二、三问数字不同，此处按原文保留）
        \end{enumerate}

  \item 单 index block：1024 指针，每数据块 4\,KB。
        $1024 \times 4\,\text{KB} = \mathbf{4\,MB}$。 \lec{14,15}

  \item 若支持多级索引：3 个 index blocks 可构成
        \begin{itemize}[nosep]
          \item 1 个一级 index block $\to$ 指向 1024 数据块；
          \item 1 个二级 index block $\to$ 指向 1024 个一级 index block
                $\to 1024^2$ 数据块；
          \item 1 个三级 index block $\to 1024^3$ 数据块。
        \end{itemize}
        极端情况下（三级间接 + 两级间接 + 一级间接），
        $\approx 1024^3 \times 4\,\text{KB} \approx \mathbf{4\,TB}$。

        若仅为 3 个平级 index block：$3 \times 1024 \times 4\,\text{KB}
        = \mathbf{12\,MB}$。

        （两种理解均可，取决于是否允许多级索引结构。） \lec{15,16}

  \item
        \begin{enumerate}[label=(\roman*)]
          \item \textbf{time of creation} $\to$ \textbf{inode}。
                创建时间是文件元数据（metadata），UNIX 系统存于 inode 中
                （如 \texttt{st\_ctime}）。
          \item \textbf{file name} $\to$ \textbf{directory}。
                Directory 的核心功能就是从人类可读的名字 $\to$ inode 号的映射。
                inode 本身不存文件名（一个文件可以有多个硬链接，即多个名字）。
          \item \textbf{access rights} $\to$ \textbf{inode}。
                权限位（r/w/x）及 owner/group 存于 inode 的 \texttt{mode} 字段。
          \end{enumerate}
        \lec{14,16}

  \item \textbf{Contiguous allocation} 性能分析：

        \begin{itemize}[nosep]
          \item \textbf{Reading}：极优。文件在磁盘上连续存放，
                顺序读只需一次 seek + 旋转延迟后可持续传输，I/O 带宽利用率最高。
          \item \textbf{Writing}：顺序写极优（追加）；随机写若需扩展文件
                则可能需整体搬迁。
          \item \textbf{Creating}：困难。需在磁盘上找到足够大的连续空闲区域
                （类似内存的 external fragmentation 问题）。
                随使用时间增长，空闲空间碎片化，大文件创建越来越难。
        \end{itemize}

        综合评价：contiguous 在读/写吞吐量上最优，但空间管理代价最高。
        现代文件系统（ext4, NTFS）通常混合使用 extent-based 分配 +
        按需碎片整理，而非纯 contiguous。 \lec{14,15}

\end{enumerate}

\end{document}
